%/*% pdflatex -file-line-error -interaction nonstopmode %

\documentclass{article}
\usepackage{graphicx}
\usepackage{amsmath}
\usepackage{amsfonts}
\usepackage{bookmark}
\usepackage{booktabs}
\usepackage{exercises}
\graphicspath{ {./images/test} }
\usepackage{indentfirst}
\usepackage{amsthm}
\usepackage{yhmath}
\usepackage{placeins}
\newcommand*{\oldneg}{\mathord{\sim}}
\newcommand\showdiv[1]{\overline{\smash{)}#1}}
\newcommand\Mydiv[1]{%
\smash{\raise.3ex\hbox{$\big)$}}\mkern-8mu
        \overline{\enspace\strut#1}}
\setlength{\tabcolsep}{5pt}
\theoremstyle{definition}
\newtheorem{definition}{Definition}

\setlength{\parindent}{.5in}
\setlength{\parskip}{12pt}
\renewcommand{\arraystretch}{2}

\title{3.2 Answers}
\begin{document}
\maketitle


\section*{1}
Like radicals are defined to be radicals with the same index and the same radicand. We can see that $2\sqrt{5}$ and $2\sqrt[3]{5}$ have different 
indeces (1 and 3), meaning they are not like radicals.

\section*{2}
A radical expression with index $n$ is in simplest form iff the radicand has no factors of index $n$ and any denominator in the radicand is 
rationalized.

\section*{12}
(I was unsure if they said a simplified rational expression has no radicals in the denominator, or only rational radicals (or if those are the same thing))
\newline
\[ \frac{64^\frac{5}{9}\times 64^\frac{2}{9}}{4^\frac{3}{4}}=\frac{64^{\frac{5}{9}+\frac{2}{9}}}{4^\frac{3}{4}} \]
\[ =\frac{64^\frac{7}{9}}{4^\frac{3}{4}}=\frac{(4^3)^\frac{7}{9}}{4^\frac{3}{4}}=\frac{4^\frac{21}{9}}{4^\frac{3}{4}} \]
\[ =\frac{4^\frac{7}{3}}{4^\frac{3}{4}}=\frac{4^\frac{28}{12}}{4^\frac{9}{12}}=4^\frac{19}{12} \]
\[ =(2^2)^\frac{19}{12}=2^\frac{38}{12}=2^\frac{19}{6}=2^3\sqrt[6]{2} \]

\section*{13}
\[ (16^\frac{5}{9}\times 5^\frac{7}{9})^{-3}=16^{-\frac{5}{3}}\times 5^{-\frac{7}{3}}=\frac{1}{16^\frac{5}{3}\times 5^\frac{7}{3}} \]
\[ \frac{1}{(4^2)^\frac{5}{3}\times 25\sqrt[3]{5}}=\frac{1}{64\sqrt[3]{4}\times 25\sqrt[3]{5}} \]
\[ =\frac{\sqrt[3]{4^2}\times \sqrt[3]{5^2}}{64\times 4\times 25\times 5}=\frac{2\sqrt[3]{2}\times \sqrt[3]{5^2}}{64\times 4\times 25\times 5} \]

\section*{14}
\[ \frac{13^\frac{3}{7}}{13^\frac{5}{7}}=13^{\frac{3}{7}-\frac{5}{7}}=13^{-\frac{2}{7}}=\frac{1}{13^\frac{2}{7}} \]

\section*{21}
\[ \frac{\sqrt[4]{36}\times \sqrt[4]{9}}{\sqrt[4]{4}}=\frac{\sqrt[4]{324}}{\sqrt[4]{4}}=\frac{\sqrt[4]{3^4\times 4}}{\sqrt[4]{4}} \]
\[ =\frac{3\sqrt[4]{4}}{\sqrt[4]{4}}=3 \]

\section*{22}
\[ \frac{\sqrt[4]{8}\times\sqrt[4]{16}}{\sqrt[8]{2}\times\sqrt[8]{3}}=\frac{2\sqrt[4]{8}}{\sqrt[8]{6}} \]
\[ =\frac{2\sqrt[4]{8}\times\sqrt[8]{6^7}}{6} \]

\section*{23}
C

\section*{30}
\[ \sqrt[6]{\frac{81}{4}}=\frac{\sqrt[6]{81}}{\sqrt[6]{4}}=\frac{\sqrt[6]{81}}{\sqrt[6]{4}}\times\frac{\sqrt[6]{2^4}}{\sqrt[6]{2^4}} \]
\[ =\frac{\sqrt[6]{1296}}{2} \]

\section*{31}
\[ \frac{\sqrt[3]{9}}{\sqrt[5]{27}}=\frac{\sqrt[3]{9}}{\sqrt[5]{27}}\times\frac{\sqrt[5]{3^2}}{\sqrt[5]{3^2}} \]
\[ =\frac{\sqrt[3]{9}\times\sqrt[5]{9}}{3}=\frac{9^\frac{8}{15}}{3}=\frac{(3^2)^\frac{8}{15}}{3}=3^\frac{1}{15} \]

\section*{39}
\[ 2\sqrt[4]{1250}-8\sqrt[4]{32}=2\sqrt[4]{1250}-16\sqrt[4]{2}=10\sqrt[4]{2}-16\sqrt[4]{2}=-6\sqrt[4]{2} \]

\section*{40}
\[ 5\sqrt[3]{48}-\sqrt[3]{750}=10\sqrt[3]{6}-5\sqrt[3]{6}=5\sqrt[3]{6} \]

\section*{41}
The error here is that you can only combine $like$ radicals, that is, radicals with the same radicand and the same index. In this problem, 
the radicals have different radicands. The problem can be corrected by not doing any steps, as the expression is already fully simplified.

\section*{49}
\[ (\sqrt[3]{x^2}\times \sqrt[6]{x^4})^{-3}=(\sqrt[3]{x^2})^{-3}\times (\sqrt[6]{x^4})^{-3} \]
\[ =\sqrt[3]{x^{-6}}\times \sqrt[6]{x^{-12}}=x^{-2}\times x^{-2}=x^{-4}=\frac{1}{x^4} \]

\section*{50}
\[ \dfrac{\sqrt[3]{x}\times\sqrt{x^5}}{\sqrt{25x^{16}}}=\frac{x^\frac{17}{6}}{5x^8} \]
\[ =\frac{x^\frac{17}{6}-x^8}{5}=\frac{1}{5x^\frac{31}{6}}=\frac{1}{5\sqrt[6]{x^{31}}} \]

\section*{51}
$x^\frac{5}{4}$ and $x^\frac{1}{2}$.

\section*{58}
\[ \sqrt{\frac{20x^3y^2}{9xz^3}}=\frac{\sqrt{20x^3y^2}}{\sqrt{9xz^3}}=\frac{2xy\sqrt{5x}}{3z\sqrt{xz}} \]

\section*{59}
\[ \frac{\sqrt[4]{x^6}}{\sqrt[7]{x^5}}=\frac{x\sqrt[4]{x^2}}{\sqrt[7]{x^5}}\times\frac{\sqrt[7]{x^2}}{\sqrt[7]{x^2}} \]
\[ =\frac{x\sqrt[4]{x^2}\sqrt[7]{x^2}}{x}=\sqrt[4]{x^2}\sqrt[7]{x^2}=x^\frac{11}{14} \]

\section*{64}
\[ x\sqrt{9x^3}-2\sqrt{x^5}=3x\sqrt{x^3}-2x\sqrt{x^3}=x\sqrt{x^3}=x^2\sqrt{x} \]

\section*{65}
\[ y\sqrt[4]{32x^6}+\sqrt[4]{162x^2y^4}=2xy\sqrt[4]{2x^2}+3y\sqrt[4]{2x^2} \]

\section*{68}
\[ A=\frac{1}{2}3x^\frac{1}{3}\times 4x^\frac{1}{3}=\frac{1}{2}\times 12x^\frac{2}{3}=6x^\frac{2}{3} \]
We know that the third side is \[ \sqrt{(3x^\frac{1}{3})^2+(4x^\frac{1}{3})^2}=\sqrt{9x^\frac{2}{3}+16x^\frac{2}{3}} \]
\[ =\sqrt{25x^\frac{2}{3}}=5x^\frac{1}{3} \] meaning \[ P=3x^\frac{1}{3}+4x^\frac{1}{3}+5x^\frac{1}{3}=12x^\frac{1}{3} \]

\section*{69}
C

\section*{76}
\[ 10x^{0.6}+(4x^{0.3})^2=10x^\frac{6}{10}+16x^\frac{6}{10}=26\sqrt[10]{x^6} \]

\section*{77}
\[ 15x^{0.3}-(2x^{0.1})^3=15x^\frac{3}{10}-8x^\frac{3}{10}=7\sqrt[10]{x^3} \] 

\section*{80}
\[ (\frac{x^\pi}{x^\frac{\pi}{3}})^2=(x^{\pi-\frac{\pi}{3}})^2=(x^\frac{2\pi}{3})^2=x^\frac{4\pi}{3} \]

\section*{81}
\[ x^2y^{\sqrt{2}}+3x^2y^{\sqrt{2}}=4x^2y^{\sqrt{2}} \]

\section*{85}
We know that $d=1.9[(5.5\times 10^{-4})l]^\frac{1}{2}$. Thus, if $l=10$, then
\[ d=1.9[(5.5\times 10^{-4}\times 10)]^\frac{1}{2}=1.9[(5.5\times 10^{-3})]^\frac{1}{2} \]
\[ \frac{1.9\sqrt{5.5}}{10\sqrt{10}} \]

\section*{88}
\[ d=\dfrac{v_0\sqrt{(v_0)^2+2g(0)}}{g}=\dfrac{v_0\sqrt{v_0^2}}{g}=\dfrac{v_0^2}{g} \]

\section*{89}
(a): If $V=\frac{4}{3}\pi r^3$ then $r=\sqrt[3]{\frac{3V}{4\pi}}$.\newline


(b): \[ S=4\pi r^2 \] \[ S=4\pi(\sqrt[3]{\frac{3V}{4\pi}})^2=4\pi\times (4\pi)^{-\frac{2}{3}}\times (3V)^\frac{2}{3} \]
\[ =(4\pi)^\frac{1}{3}\times (3V)^\frac{2}{3} \]

(c): The surface area of the first ballon is $(4\pi)^\frac{1}{3}\times (3x)^\frac{2}{3}$ while the surface area of the second balloon is 
$ (4\pi)^\frac{1}{3}\times (6x)^\frac{2}{3}$, meaning the ratio of the first to the second balloon is:
\[ \frac{(4\pi)^\frac{1}{3}\times (3x)^\frac{2}{3}}{(4\pi)^\frac{1}{3}\times (6x)^\frac{2}{3}}= \]
\[ =\frac{(6x)^\frac{2}{3}}{(3x)^\frac{2}{3}}=\frac{6x}{3x}^\frac{2}{3}=2^\frac{2}{3} \]

\section*{90}
(I didn't quite understand what was being asked here.)

\end{document}

